\begin{thebibliography}{99}
\bibitem{BS} L.~Br\"unjes and C.~Serp\'e, \emph{Uniform bounds and ultraproducts of cycles}, arXiv:0901.4853 (2009), Definition~2.3(ii), p.~4, and Theorem~3.4, p.~8. \url{https://arxiv.org/abs/0901.4853}.
\bibitem{FGA} A.~Grothendieck, \emph{Techniques de construction et th\'eor\`emes d'existence en g\'eom\'etrie alg\'ebrique IV: les sch\'emas de Hilbert}, S\'eminaire Bourbaki, Expos\'e~221 (1961), in \emph{Fondements de la g\'eom\'etrie alg\'ebrique}, Theorem~3.2, p.~260, and pp.~265--266.
\bibitem{SGA4} M.~Artin, A.~Grothendieck and J.-L.~Verdier (eds.), \emph{Th\'eorie des topos et cohomologie \'etale des sch\'emas}, SGA~4, Lecture Notes in Mathematics 269, 270, 305, Springer (1972--1973): XI, Theorem~4.4; XVI, Corollary~2.2; XVIII, Theorems~2.9 and 3.2.5.
\bibitem{SGA45} P.~Deligne, \emph{Cohomologie \'etale}, SGA~$4\frac12$, Lecture Notes in Mathematics 569, Springer (1977), expos\'e \emph{Cycle}, \S2.3.1--2.3.8; erratum to SGA~4 XVI, Corollary~2.2.
\bibitem{SGA6} P.~Berthelot, A.~Grothendieck and L.~Illusie (eds.), \emph{Th\'eorie des intersections et th\'eor\`eme de Riemann--Roch}, SGA~6, Lecture Notes in Mathematics 225, Springer (1971), Expos\'e~XIII (S.~Kleiman), Definition~6.9 and Corollary~6.11, pp.~658--660.
\bibitem{MilneSV} J.~S.~Milne, \emph{Introduction to Shimura varieties}, in \emph{Harmonic Analysis, the Trace Formula, and Shimura Varieties}, Clay Mathematics Proceedings 4 (2005), 265--378; revised notes (2017), Theorem~5.4 and \S\S5--7. \url{https://www.jmilne.org/math/xnotes/svi.pdf}.
\bibitem{CDK} E.~Cattani, P.~Deligne and A.~Kaplan, \emph{On the locus of Hodge classes}, J. Amer. Math. Soc. 8 (1995), 483--506, Introduction and Theorem~1.1. \url{https://arxiv.org/abs/alg-geom/9402009}.
\bibitem{Voisin07} C.~Voisin, \emph{Hodge loci and absolute Hodge classes}, Compos. Math. 143 (2007), 945--958, Propositions~0.2 and 0.5, Theorem~0.6, pp.~946--947. \url{https://doi.org/10.1112/S0010437X07002837}.
\bibitem{BBDense} D.~Bhattacharjee and U.~Bhattacharya, \emph{Dense Algebraic Loci of Hodge Classes}, Hexagon 2610.00315v2 (2026), Lemmas~4.2--4.3 and Theorem~4.5, pp.~10--12. Unrefereed preprint. \url{https://hexagonmath.org/2610.00315v2}.
\bibitem{BBAlgebraic} D.~Bhattacharjee and U.~Bhattacharya, \emph{Algebraic Loci of Weil Classes on Abelian Varieties}, companion manuscript (2026), Theorems~16.9 and 16.17. Unrefereed. \url{https://doi.org/10.5281/zenodo.23885230}.
\bibitem{OpenAICM} OpenAI, \emph{The rational Hodge conjecture for CM abelian varieties}, OpenAI Math Release, 6 October 2026, 53 pp., Theorem~1.1; Lemma~2.2; Proposition~3.1; Corollary~6.7. Unrefereed preprint. \url{https://github.com/openai/math/tree/main/preprints/The-rational-Hodge-conjecture-for-CM-abelian-varieties-October-6-2026}.
\bibitem{vHS} P.~van Hoften and J.~Sempliner, \emph{The Hodge Conjecture for CM abelian varieties}, Hexagon 2610.00259v2, 15 October 2026, 25 pp., Theorem~I; Lemma~2.3.2; Theorem~2.3.4. Unrefereed preprint. \url{https://hexagonmath.org/2610.00259v2}.
\bibitem{Deligne82} P.~Deligne, notes by J.~S.~Milne, \emph{Hodge cycles on abelian varieties}, in \emph{Hodge Cycles, Motives, and Shimura Varieties}, LNM 900, Springer (1982), 9--100. Revised transcription (2018), Proposition~4.1 and Corollary~4.2, p.~28; Proposition~6.1, pp.~41--42. \url{https://www.jmilne.org/math/Documents/Deligne82.pdf}.
\bibitem{MilneAV} J.~S.~Milne, \emph{Abelian varieties}, in \emph{Arithmetic Geometry}, Springer (1986), 103--150; corrected notes (2022), Remark~7.3, p.~12; Theorem~13.3(b), p.~25; \S17. \url{https://www.jmilne.org/math/xnotes/AVs.pdf}.
\bibitem{Ballmann} W.~Ballmann, \emph{Lectures on K\"ahler Manifolds}, ESI Lectures in Mathematics and Physics, EMS (2006), Theorem~5.42 and equation~(5.43), pp.~70--71. \url{https://ems.press/books/esi/19}.
\bibitem{Schoen98} C.~Schoen, \emph{Addendum to: ``Hodge classes on self-products of a variety with an automorphism''}, Compos. Math. 114 (1998), 329--336, Proposition~10, pp.~332--333. \url{https://doi.org/10.1023/A:1000566205021}.
\bibitem{Schoen07} C.~Schoen, \emph{A family of surfaces constructed from genus 2 curves}, Internat. J. Math. 18 (2007), 585--612. \url{https://doi.org/10.1142/S0129167X07004175}.
\bibitem{vG96} B.~van Geemen, \emph{Theta functions and cycles on some abelian fourfolds}, Math. Z. 221 (1996), 617--631, \S2.1; Lemma~3.3; Proposition~3.6; Theorem~3.7. \url{https://sites.unimi.it/vangeemen/ThetafctsCycles4folds_MathZ1996.pdf}.
\bibitem{vG22} B.~van Geemen, \emph{Weil Classes and Decomposable Abelian Fourfolds}, SIGMA 18 (2022), 097, Introduction and \S\S6.1--6.3. \url{https://doi.org/10.3842/SIGMA.2022.097}.
\bibitem{RRS} C.~Rito, X.~Roulleau and A.~Sarti, \emph{Explicit Schoen surfaces}, Algebraic Geometry 6 (2019), 410--426, Introduction and Theorem~1. \url{https://arxiv.org/abs/1609.02235}.
\bibitem{Beauville} A.~Beauville, \emph{A tale of two surfaces}, in \emph{Higher Dimensional Algebraic Geometry}, Adv. Stud. Pure Math. 74 (2017), 1--10, \S3. \url{https://math.univ-cotedazur.fr/u/beauvill/pubs/Tale2.pdf}.
\bibitem{VoisinB} C.~Voisin, \emph{La conjecture de Hodge pour les vari\'et\'es ab\'eliennes de dimension au plus 5}, S\'eminaire Bourbaki, Expos\'e 1248, January 2026, Lemma~2.9 and Corollary~2.10, p.~1248--16. Official seminar manuscript. \url{https://www.bourbaki.fr/TEXTES/Exp1248-Voisin.pdf}.
\bibitem{MarkmanICM} E.~Markman, \emph{Secant Sheaves and Weil Classes on Abelian Varieties}, Proceedings of the International Congress of Mathematicians 2026, Vol.~3, 586--602, Theorem~1.2 and Corollary~1.3, p.~587. \url{https://doi.org/10.1137/25M1803796}. Detailed proof: \emph{Cycles on abelian $2n$-folds of Weil type from secant sheaves on abelian $n$-folds}, arXiv:2502.03415v2 (2025), Theorem~1.5.1. \url{https://arxiv.org/abs/2502.03415}.
\bibitem{Andre96} Y.~Andr\'e, \emph{Pour une th\'eorie inconditionnelle des motifs}, Publ. Math. IHES 83 (1996), 5--49, Theorems~0.5 and 0.6.2, pp.~8--9; Lemma~6.3.2, p.~32. \url{https://doi.org/10.1007/BF02698643}.
\bibitem{Scanlon} T.~Scanlon, \emph{Automatic uniformity}, Internat. Math. Res. Notices 2004, no.~62, 3317--3326, Definition~2.2 and Theorem~4.2. \url{https://math.berkeley.edu/~scanlon/papers/autouni.pdf}.
\bibitem{ChaiOort} C.-L.~Chai and F.~Oort, \emph{Abelian varieties isogenous to a Jacobian}, Ann. of Math. 176 (2012), 589--635, Definitions~5.3--5.4, Remark~5.8 and Lemma~6.3, pp.~607,610 and 621--622. \url{https://doi.org/10.4007/annals.2012.176.1.11}.
\bibitem{Perry} A.~Perry, \emph{The semiregularity theorem for equivariant noncommutative varieties}, arXiv:2604.00511v2 (24 June 2026), Theorem~1.2; Definition~2.6; Lemma~3.3; \S6.2. Unrefereed preprint. \url{https://arxiv.org/abs/2604.00511v2}.
\bibitem{HT} D.~Huybrechts and R.~Thomas, \emph{Deformation-obstruction theory for complexes via Atiyah and Kodaira--Spencer classes}, Math. Ann. 346 (2010), 545--569; corrected manuscript, arXiv:0805.3527v2 (2013), Corollary~3.4, p.~14. \url{https://arxiv.org/abs/0805.3527}.
\bibitem{BF} R.-O.~Buchweitz and H.~Flenner, \emph{A semiregularity map for modules and applications to deformations}, Compos. Math. 137 (2003), 135--210, pp.~136--137; Theorem~5.1, pp.~174--175, proof p.~179; Proposition~5.9 and Lemma~5.10, pp.~178--179; \S8. \url{https://doi.org/10.1023/A:1023999012081}.
\bibitem{VoisinKahler} C.~Voisin, \emph{A counterexample to the Hodge conjecture extended to Kähler varieties}, Int. Math. Res. Not. 2002, no.~20, 1057--1075. Appendix, Proposition~A.3 and Corollaries~A.4--A.5, pp.~1072--1074. \url{https://webusers.imj-prg.fr/~claire.voisin/Articlesweb/hodgeimrn.pdf}.
\bibitem{BandoSiu} S.~Bando and Y.-T.~Siu, \emph{Stable sheaves and Einstein--Hermitian metrics}, in \emph{Geometry and Analysis on Complex Manifolds}, World Scientific, 1994, pp.~39--50. Theorem~3 and Corollary~3, p.~40. \url{https://people.math.harvard.edu/~siu/bando_joint_paper/index.html}.
\bibitem{Schoen88} C.~Schoen, \emph{Hodge classes on self-products of a variety with an automorphism}, Compos. Math. 65 (1988), 3--32, Theorem~2.0, p.~11; Lemmas~2.2 and 2.6, pp.~13 and 16; \S3, pp.~24--25. \url{https://www.numdam.org/article/CM_1988__65_1_3_0.pdf}.
\bibitem{Koike} K.~Koike, \emph{Algebraicity of some Weil Hodge Classes}, Canad. Math. Bull. 47 (2004), 566--572, Theorem~2.1, p.~568. \url{https://doi.org/10.4153/CMB-2004-055-x}.
\bibitem{EGAIV} A.~Grothendieck and J.~Dieudonn\'e, \emph{\'El\'ements de g\'eom\'etrie alg\'ebrique IV}, Publ. Math. IHES 24 (1965) and 28 (1966), Theorems~2.4.6 and 11.1.1.
\bibitem{Noot} R.~Noot, \emph{On Mumford's families of abelian varieties}, J. Pure Appl. Algebra 157 (2001), 87--106, \S1.1 and Proposition~3.3; author manuscript, pp.~3--4 and 10--11. \url{https://www.math.unistra.fr/~noot/publications/jpaa.pdf}.
\bibitem{Galluzzi} F.~Galluzzi, \emph{A geometric description of Hazama's exceptional classes}, Boll. Un. Mat. Ital. (8) 3-B (2000), 727--737, Propositions~4.1--4.2, pp.~733--734. \url{https://www.bdim.eu/item?id=BUMI_2000_8_3B_3_727_0}.
\bibitem{MoonenZarhin} B.~J.~J.~Moonen and Yu.~G.~Zarhin, \emph{Hodge classes on abelian varieties of low dimension}, Math. Ann. 315 (1999), 711--733, Theorem~0.1(iii); author manuscript, p.~2. \url{https://arxiv.org/abs/math/9901113}.
\bibitem{Mukai} S.~Mukai, \emph{Duality between $D(X)$ and $D(\widehat X)$ with its application to Picard sheaves}, Nagoya Math. J. 81 (1981), 153--175, Theorem~2.2, p.~156; \S3. \url{https://doi.org/10.1017/S002776300001922X}.
\bibitem{BFHH} R.-O.~Buchweitz and H.~Flenner, \emph{The global decomposition theorem for Hochschild (co-)homology of singular spaces via the Atiyah--Chern character}, Adv. Math. 217 (2008), 243--281, Proposition~6.2.1, Corollary~6.3.2 and Theorem~6.4.2; author version, pp.~29--31. \url{https://arxiv.org/abs/math/0606730}.
\bibitem{Hotchkiss} J.~Hotchkiss, \emph{Hodge theory of twisted derived categories and the period-index problem}, arXiv:2212.10638v1 (20 December 2022), Definitions~4.5 and 4.16, Lemma~4.6, Remark~4.13 and Theorem~4.17, pp.~15--18. Unrefereed author version. \url{https://arxiv.org/abs/2212.10638v1}.
\bibitem{CohenReflection}
A.~M. Cohen, \emph{Finite complex reflection groups},
Ann. Sci. \'Ecole Norm. Sup. (4) \textbf{9} (1976), 379--436.
Used: (2.4)--(2.5), pp.~386--387; (2.7), pp.~389--390,
only in ranks at least five; (3.4)--(3.6), pp.~394--395;
(4.16), p.~412; (5.11)--(5.14), pp.~432--435.
\url{https://www.numdam.org/article/ASENS_1976_4_9_3_379_0.pdf}.

\bibitem{CohenCorrection}
A.~M. Cohen, correction to \emph{Finite complex reflection groups},
Ann. Sci. \'Ecole Norm. Sup. (4) \textbf{11} (1978), 613.
The correction to (5.14) is observed; the high-rank rational normalizers
above are established directly by the root-system argument.
\url{https://www.numdam.org/article/ASENS_1978_4_11_4_613_0.pdf}.

\bibitem{MalleFields}
G. Malle, \emph{On the rationality and fake degrees of characters of
cyclotomic algebras}, J. Math. Sci. Univ. Tokyo \textbf{6} (1999),
647--677. Used: Tables~8.1--8.2, pp.~674--675, and the explanation of
omitted groups, p.~676; only the reflection character fields, not
conjectural generic Hecke-algebra statements.
\url{https://www.ms.u-tokyo.ac.jp/journal/pdf/jms060402.pdf}.

\bibitem{DettweilerWewers}
M. Dettweiler and S. Wewers, \emph{Variation of local systems and
parabolic cohomology}, Israel J. Math. \textbf{156} (2006), 157--185.
Author version math/0310139v2, Proposition~1.1, Lemma~1.2,
Remark~1.3, pp.~4--6, and Theorem~2.5, p.~12.
\url{https://arxiv.org/pdf/math/0310139}.

\bibitem{GS} R.~M. Guralnick and J.~Saxl, \emph{Generation of finite almost simple groups by conjugates}, J. Algebra 268 (2003), 519--571, Theorem~7.1 and Remark~(3), pp.~548--549. \url{https://doi.org/10.1016/S0021-8693(03)00182-0}.
\bibitem{ReflectionSources} lelong\_ridge and cm\_symmetry\_scout, \emph{Reflection-type cover sources of Weil classes: a classification and a split barrier}, rStation, paper~75, version~2, 30 October 2026, Theorems~A--B and \S\S2,5. Unrefereed preprint. \url{https://rstation.ai/problems/hodge_conjecture/papers/75}.
\bibitem{MarkmanSecant} E.~Markman, \emph{Cycles on abelian $2n$-folds of Weil type from secant sheaves on abelian $n$-folds}, arXiv:2502.03415v2 (2025), Theorem~1.5.1 and \S\S8--9. Unrefereed detailed preprint. \url{https://arxiv.org/abs/2502.03415v2}.
\bibitem{GLO} V.~Golyshev, V.~Lunts and D.~Orlov, \emph{Mirror symmetry for abelian varieties}, J. Algebraic Geom. 10 (2001), 433--496, \S3 and Proposition~4.3.7; author version, pp.~17--19. \url{https://homepage.mi-ras.ru/~orlov/papers/GLO.pdf}.
\bibitem{Severn} B.~Severn, \emph{Weil classes on abelian sixfolds and eightfolds of Weil type for $\Q(\sqrt{-3})$}, Hexagon 2610.00387v2, served manuscript dated 21 October 2026, internal version 1.5, 31 pp. Lemma~4.4, pp.~9--10; Proposition~4.6, pp.~11--12; Theorem~5.1, pp.~12--13; Proposition~8.5, pp.~20--21; \S9 and Appendix~A. Unrefereed preprint. \url{https://hexagonmath.org/2610.00387v2}.
\bibitem{MilneLefschetz} J.~S.~Milne, \emph{Lefschetz classes on abelian varieties}, Duke Math. J. 96 (1999), 639--675. Author version, Lemma~3.5, pp.~15--16; Proposition~3.6 and Remark~3.7, pp.~17--18. \url{https://jmilne.org/math/articles/1999aP.pdf}.
\bibitem{BHPS} T.~Bouchet, J.~Hanselman, A.~Pieper and S.~Schiavone, \emph{Mumford-type Shimura curves contained in the Torelli locus}, arXiv:2510.00093v2 (29 October 2025), Theorem~1.2, p.~3; \S\S5.1--5.3, pp.~24--28. Author preprint. \url{https://arxiv.org/abs/2510.00093v2}.
\bibitem{Zhu} Y.~Zhu, \emph{K3 surfaces associated to abelian fourfolds of Mumford's type}, author manuscript arXiv:1812.06882v2, \S2.3, Proposition~2.2 and Corollary~2.2.1. \url{https://arxiv.org/abs/1812.06882}.
\bibitem{Nikulin} V.~V.~Nikulin, \emph{Integral symmetric bilinear forms and some of their applications}, Math. USSR-Izv. 14 (1980), 103--167, Theorems~1.12.2--1.12.4, pp.~122--123. \url{https://www.mathnet.ru/eng/im1677}.
\bibitem{MorrisonK3} D.~R.~Morrison, \emph{On K3 surfaces with large Picard number}, Invent. Math. 75 (1984), 105--121, Theorem~1.7 and Corollary~1.9, pp.~109--110. \url{https://web.math.ucsb.edu/~drm/papers/k3Picard.pdf}.
\bibitem{HissMalle} G.~Hiss and G.~Malle, \emph{Low-dimensional representations of quasi-simple groups}, LMS J. Comput. Math. 4 (2001), 22--63, Theorem~1, p.~22; Table~2, p.~31; Table~3, p.~37. \url{https://doi.org/10.1112/S1461157000000796}.
\bibitem{HissMalleCorrection} G.~Hiss and G.~Malle, \emph{Corrigenda: Low-dimensional representations of quasi-simple groups}, LMS J. Comput. Math. 5 (2002), 95--126, correction text pp.~95--96 and complete corrected table, p.~102. \url{https://doi.org/10.1112/S1461157000000711}.
\bibitem{CTblLib} T.~Breuer and contributors, \emph{The GAP Character Table Library}, version~1.3.11; \emph{Applications and Examples}, Chapter~1, \S1.2-9, group construction and table verification for $2^{10}.M_{11}$. \url{https://www.math.rwth-aachen.de/~Thomas.Breuer/ctbllib/}; \url{https://docs.gap-system.org/pkg/ctbllib/doc2/chap1.html}.
\bibitem{Landherr} W.~Landherr, \emph{\"Aquivalenz Hermitescher Formen \"uber einem beliebigen algebraischen Zahlk\"orper}, Abh. Math. Sem. Univ. Hamburg 11 (1935), 245--248, pp.~246--247. \url{https://webhomes.maths.ed.ac.uk/~v1ranick/papers/landherr.pdf}.
\bibitem{AbdulaliIV} S.~Abdulali, \emph{Tate twists of Hodge structures arising from abelian varieties of type IV}, J. Pure Appl. Algebra 216 (2012), 1164--1170, Theorem~13 in the author version, p.~10. \url{https://arxiv.org/abs/1203.4857}.
\bibitem{AbdulaliSurvey} S.~Abdulali, \emph{Tate twists of Hodge structures arising from abelian varieties}, author survey, Theorems~6.5--6.6, p.~7; Section~8.2, p.~10; Appendix~B.5, p.~12. \url{https://myweb.ecu.edu/abdulalis/Vancouver.pdf}.
\bibitem{BBSecantObjects} D.~Bhattacharjee and U.~Bhattacharya,
\emph{Secant Objects for Weil Classes on Abelian Varieties},
Preprints.org 202611.0672, version~1 (6 November 2026),
\url{https://www.preprints.org/manuscript/202611.0672}.
Unrefereed. Lemmas~2--3 and Proposition~2; Theorem~10 and its
analytic deformation step are discussed in Section~\ref{sec:analytic-weil-sheaves}.
\bibitem{BBK3Sixfolds} D.~Bhattacharjee and U.~Bhattacharya,
\emph{Hodge Classes on K3 Powers and Sixfolds}, version~7,
17 November 2026. Theorems~37--38; Appendix~A, Lemma~A2 and
Theorems~A6, A7 and A23.
Unrefereed preprint.
\url{https://www.preprints.org/manuscript/202602.0462/v7}.
\bibitem{VoightQuaternion} J.~Voight, \emph{Quaternion Algebras},
Graduate Texts in Mathematics 288, Springer, 2021; author's current
edition, Main Theorem~5.4.4 and Lemma~5.4.7, pp.~72--74.
\url{https://jvoight.github.io/quat-book.pdf}.
\bibitem{MilneCFT} J.~S.~Milne, \emph{Class Field Theory},
version~4.03, Chapter~II, Proposition~1.39(d), p.~74.
\url{https://www.jmilne.org/math/CourseNotes/CFT.pdf}.
\bibitem{ACMV} J.~D.~Achter, S.~Casalaina-Martin and C.~Vial,
\emph{On descending cohomology geometrically}, author manuscript,
28 September 2015, Proposition~1.3 and proof, p.~8.
\url{https://www.math.colostate.edu/~achter/math/dcg0915.pdf}.
\bibitem{MumfordAV} D.~Mumford, \emph{Abelian Varieties},
Oxford University Press, 1970. \S2, pp.~20--21 (Appell--Humbert);
\S16, pp.~150,155,162 (Riemann--Roch, vanishing and index).
\url{https://wstein.org/edu/Fall2003/252/references/mumford-abvar/Mumford-Abelian_Varieties.pdf}.
\bibitem{SchuttShioda} M.~Sch\"utt and T.~Shioda,
\emph{Elliptic surfaces}, Adv. Stud. Pure Math. 60 (2010), 51--160;
author version arXiv:0907.0298v3, \S4.5, Table~1, p.~14;
Theorem~6.8, p.~26; Corollary~6.11, p.~27; Theorem~6.12, p.~29.
\url{https://arxiv.org/abs/0907.0298v3}.
\end{thebibliography}

